%% Question bank chapter file - generated by bank_tools.py build (do not edit while the bank is still being built:
%% it is regenerated from the per-test data in Question Bank\data). Every question is one bankq block:
%%   \begin{bankq}[id=..., source={...}, paper=.., qno=.., type=..., subtopic={...}, answer={...}, answer_src=key|solved,
%%                 working={...}, difficulty=1-10, diff_parts={K S I C F}, diff_note={...},
%%                 group=<passage id or empty>, also_in={duplicate copies in other tests}]
%% difficulty = min(10, K+S+I+C+F-1): K concepts 1-3, S steps 1-4, I insight 0-2, C computation 0-2, F format 0-1
%%   ...question LaTeX, options, figures (\bfig / \bfigi -> img/)...
%%   \end{bankq}
%% Passages shared by several questions are bankpassage blocks; their questions carry group=<passage id>.
\documentclass[10pt]{article}
\usepackage[a4paper,margin=1.8cm]{geometry}
\usepackage[T1]{fontenc}
\usepackage{amsmath,amssymb,mathtools}
\usepackage{graphicx}
\usepackage[export]{adjustbox}
\usepackage{xcolor,needspace,keyval}
\usepackage{enumitem}
\usepackage{booktabs,array,tabularx}
\usepackage[hidelinks]{hyperref}
\setlength{\parindent}{0pt}
\setlength{\parskip}{3pt}
\makeatletter
\define@key{bq}{id}{\def\bq@id{#1}}\define@key{bq}{source}{\def\bq@src{#1}}\define@key{bq}{paper}{\def\bq@paper{#1}}
\define@key{bq}{qno}{\def\bq@qno{#1}}\define@key{bq}{type}{\def\bq@type{#1}}\define@key{bq}{subtopic}{\def\bq@sub{#1}}
\define@key{bq}{answer}{\def\bq@ans{#1}}\define@key{bq}{answer_src}{\def\bq@asrc{#1}}\define@key{bq}{working}{\def\bq@work{#1}}
\define@key{bq}{group}{\def\bq@group{#1}}\define@key{bq}{also_in}{\def\bq@also{#1}}
\define@key{bq}{difficulty}{\def\bq@diff{#1}}\define@key{bq}{diff_parts}{\def\bq@dparts{#1}}\define@key{bq}{diff_note}{\def\bq@dnote{#1}}
\newenvironment{bankq}[1][]{\def\bq@work{}\def\bq@also{}\def\bq@group{}\def\bq@diff{?}\def\bq@dparts{}\def\bq@dnote{}\setkeys{bq}{#1}%
  \par\needspace{4\baselineskip}\medskip\noindent{\small\textbf{\bq@id}\hfill\color{gray}[\bq@src, Paper \bq@paper, Q\bq@qno]
  \bq@type\ $\cdot$ \bq@sub\ $\cdot$ \textbf{D\bq@diff}}\par\nopagebreak\relax}
  {\par{\small\color{teal}Answer: \textbf{\bq@ans}\ifx\bq@work\@empty\else\ (solved: \bq@work)\fi
  \ifx\bq@also\@empty\else\par\color{gray}Also in: \bq@also\fi}\par}
\define@key{bp}{id}{\def\bp@id{#1}}\define@key{bp}{source}{\def\bp@src{#1}}\define@key{bp}{questions}{\def\bp@qs{#1}}
\newenvironment{bankpassage}[1][]{\setkeys{bp}{#1}\par\needspace{4\baselineskip}\bigskip\noindent
  {\small\textbf{Passage \bp@id}\hfill\color{gray}[\bp@src; questions \bp@qs]}\par\nopagebreak\relax\itshape}{\par}
\makeatother
\newcommand{\bfig}[2]{\par\smallskip{\centering\includegraphics[width=#1pt,max width=0.62\linewidth,max height=0.32\textheight]{img/#2}\par}\smallskip}
\newcommand{\bfigi}[2]{\includegraphics[width=#1pt,max width=\linewidth,max height=0.32\textheight,valign=c]{img/#2}}
\newcommand{\bmissing}[1]{\textit{[figure missing in the source conversion: \detokenize{#1}]}}
\newenvironment{qoptions}{\begin{itemize}[leftmargin=1.8em,itemsep=1pt,topsep=2pt,label={}]}{\end{itemize}}
\newcommand{\opt}[2]{\item[(#1)] #2}
\providecommand{\fwstart}{}\providecommand{\fwend}{}\providecommand{\qfullsplit}{}
\newcolumntype{L}[1]{>{\raggedright\arraybackslash}p{#1}}
\newcolumntype{C}[1]{>{\centering\arraybackslash}p{#1}}
\newcolumntype{R}[1]{>{\raggedleft\arraybackslash}p{#1}}
\providecommand{\sectionhdr}[1]{}\providecommand{\subsectionhdr}[1]{}


\begin{document}
{\Large\bfseries Mole Concept --- Question Bank}\par\smallskip
{\small Chemistry chapter C01. Questions from all converted test papers, grouped by source series. Answers: printed key where it matches the paper, otherwise solved (answer\_src=solved, with working). Exact repeats across papers appear once, with also\_in listing the other copies.}\par\medskip\hrule

\section*{Score-I papers}
\begin{bankq}[id=SI12-P1-CHEM-S2-Q04, source={Score-I - Test 12 (09 Mar 2022)}, paper=1, qno=4, type=numerical, subtopic={Concentration terms}, answer={0.2}, answer_src=solved, working={$[\mathrm{H^+}]=10^{-5}$ M so $[\mathrm{Ca^{2+}}]=5\times10^{-6}$ M $=0.2$ mg/L $=0.2$ ppm}, difficulty=2, diff_parts={K1 S1 I0 C0 F1}, diff_note={concepts: hardness in ppm from H+ released | steps: H+ concentration from pH; Ca2+ is half of it and converts to mg per litre}]
RH$_2$ (ion exchange resin) can replace Ca$^{2+}$ ions in hard water as RH$_2$ + Ca$^{2+}\to$ RCa + 2H$^+$. If 8 Litre hard water after passing through RH$_2$ has pH = 5, then hardness in parts per million of Ca$^{2+}$ is.
\end{bankq}
\begin{bankq}[id=SI12-P1-CHEM-S2-Q05, source={Score-I - Test 12 (09 Mar 2022)}, paper=1, qno=5, type=numerical, subtopic={Stoichiometry}, answer={5}, answer_src=solved, working={Each reaction removes one volume of gas per volume of H$_2$ used; all 5 L of H$_2$ react (N$_2$ in excess), so the contraction is 5 L}, difficulty=3, diff_parts={K1 S1 I1 C0 F1}, diff_note={concepts: stoichiometry of gas volume change | steps: N2 is in excess so all 5 L H2 react and each litre removes one litre of gas | I: one volume lost per volume H2 in both reactions}]
N$_2$(g) reacts with H$_2$(g) in either of the following ways depending upon \& supply of H$_2$(g).\\
N$_2$(g) + H$_2$(g) $\to$ N$_2$H$_2$(g)\\
N$_2$(g) + 2H$_2$(g) $\to$ N$_2$H$_4$(g)\\
If 3 litre N$_{2(g)}$ and 5 litre H$_{2(g)}$ are taken initially (at same temperature and pressure). Calculate the contraction in volume after the reaction (in L)
\end{bankq}

\section*{TOAS papers}
\begin{bankq}[id=TOAS01-P2-CHEM-S1-Q05, source={TOAS - Test 01 (30 Jan 2022)}, paper=2, qno=5, type=single, subtopic={Concentration terms}, answer={D}, answer_src=key, difficulty=1, diff_parts={K1 S1 I0 C0 F0}, diff_note={concepts: concentration terms | steps: molarity and percentw/v; percentw/w and molality}]
120 g glucose $[\mathrm{C_6H_{12}O_6(s)}]$ is dissolved to prepare 1 L aqueous solution of density 1.2g/mL. Identify incorrect option for the solution
\begin{qoptions}
\opt{A}{Molarity $=2/3$ M}
\opt{B}{\% w/w $=10\%$}
\opt{C}{\% w/v $=12\%$}
\opt{D}{Molality $=1.85$ mol/kg}
\end{qoptions}
\end{bankq}
\begin{bankq}[id=TOAS01-P2-CHEM-S2-Q01, source={TOAS - Test 01 (30 Jan 2022)}, paper=2, qno=1, type=numerical, subtopic={Concentration terms}, answer={1.00}, answer_src=solved, working={Cations: Na$^+$ 0.4+0.3, Ca$^{2+}$ 0.2 = 0.9 mol; anions: SO$_4^{2-}$ 0.2, Cl$^-$ 0.3+0.4 = 0.9 mol; same volume so ratio $=1$}, difficulty=2, diff_parts={K1 S1 I0 C0 F1}, diff_note={concepts: molarity of ions | steps: total moles of cations and anions; same volume so the ratio is 1}]
2M of 100 ml $\mathrm{Na_2SO_4}$ is mixed with 3M of 100 ml NaCl solution and 1M of 200 ml $\mathrm{CaCl_2}$ solution. Then find the ratio of the molarity of all cation and all anions in final solution ?
\end{bankq}
\begin{bankq}[id=TOAS01-P2-CHEM-S2-Q03, source={TOAS - Test 01 (30 Jan 2022)}, paper=2, qno=3, type=numerical, subtopic={Concentration terms}, answer={13.89}, answer_src=solved, working={28 volume: 1 L gives 28 L O$_2$ = 1.25 mol, so 2.5 mol H$_2$O$_2$ = 85 g per L; solvent = 265-85 = 180 g; molality $=2.5/0.18=13.89$}, difficulty=3, diff_parts={K1 S2 I0 C0 F1}, diff_note={concepts: volume strength and molality | steps: molarity from 28 volume; mass of H2O2 per litre; mass of solvent; molality}]
A sample of $\mathrm{H_2O_2}$ is labelled `28 volumes' at 1 atm \& 273K. Find molality of the solution if its density is 265 g/L ?
\end{bankq}
\begin{bankq}[id=TOAS02-P1-CHEM-S1-Q01, source={TOAS - Test 02 (13 Feb 2022)}, paper=1, qno=1, type=single, subtopic={Concentration terms}, answer={C}, answer_src=key, difficulty=1, diff_parts={K1 S1 I0 C0 F0}, diff_note={concepts: volume contraction on mixing | steps: total mass and final volume from the density; percentage contraction}]
On mixing 15.0 ml of ethyl alcohol of density 0.792 g $\mathrm{ml^{-1}}$ with 15 ml of pure water at $4^\circ$C, the resulting solution is found to have a density of 0.924 g $\mathrm{ml^{-1}}$. The percentage contraction in volume is :
\begin{qoptions}
\opt{A}{8 \%}
\opt{B}{2 \%}
\opt{C}{3 \%}
\opt{D}{4 \%}
\end{qoptions}
\end{bankq}
\begin{bankq}[id=TOAS02-P1-CHEM-S3-Q02, source={TOAS - Test 02 (13 Feb 2022)}, paper=1, qno=2, type=numerical, subtopic={Stoichiometry}, answer={53.25}, answer_src=key, difficulty=2, diff_parts={K1 S1 I0 C0 F1}, diff_note={concepts: empirical formula from mass data | steps: slope of the plot gives the mass of M per Cl; atomic weight}]
A metal chloride of empirical formula MCl is formed (where M is metal). In an experiment weight of metal chloride has following variation according to weight of chlorine.
\bfig{113.7}{TOAS02_p1_chem_s3_q02_1.png}
Then atomic weight of metal (M) will be :
\end{bankq}
\begin{bankq}[id=TOAS05-P1-CHEM-S1-Q05, source={TOAS - Test 05 (03 Apr 2022)}, paper=1, qno=5, type=multi, subtopic={Stoichiometry}, answer={B,D}, answer_src=key, difficulty=2, diff_parts={K1 S1 I0 C0 F1}, diff_note={concepts: empirical and molecular formula | steps: C:H ratio gives CH2; molar mass from the mass of one molecule gives C4H8}]
Hydrocarbon contains $\dfrac{600}{7}$\% by mass carbon and its one molecule weighs $93.33\times10^{-24}$ g. Which of the following statement(s) is/are correct\ $[N_A=6.0\times10^{23}]$
\begin{qoptions}
\opt{A}{Empirical formula is $\mathrm{CH}$}
\opt{B}{Empirical formula is $\mathrm{CH_2}$}
\opt{C}{Molecular formula of $\mathrm{C_3H_8}$}
\opt{D}{Molecular formula is $\mathrm{C_4H_8}$}
\end{qoptions}
\end{bankq}
\begin{bankq}[id=TOAS05-P2-CHEM-S1-Q11, source={TOAS - Test 05 (03 Apr 2022)}, paper=2, qno=11, type=single, subtopic={Stoichiometry}, answer={A}, answer_src=key, difficulty=3, diff_parts={K1 S2 I1 C0 F0}, diff_note={concepts: composition of oleum | steps: free SO3 from the label; masses in 178 g; moles of H2SO4 and free SO3; NaOH | I: oleum as H2SO4 plus free SO3}]
A sample of oleum is labelled as $\left(100+\dfrac{900}{89}\right)\%$. Find the moles of NaOH required to completely react with 178 gm of the oleum.
\begin{qoptions}
\opt{A}{4}
\opt{B}{6}
\opt{C}{3}
\opt{D}{7}
\end{qoptions}
\end{bankq}
\begin{bankq}[id=TOAS05-P2-CHEM-S1-Q12, source={TOAS - Test 05 (03 Apr 2022)}, paper=2, qno=12, type=single, subtopic={Stoichiometry}, answer={B}, answer_src=key, difficulty=1, diff_parts={K1 S1 I0 C0 F0}, diff_note={concepts: oleum with water and carbonate | steps: moles of CO2 from 1.06 g Na2CO3; volume at 300 K}]
If excess water is added into a bottle sample labelled as ``112\% oleum'' and is reacted with 1.06 gm $\mathrm{Na_2CO_3}$, then find the volume of $\mathrm{CO_2}$ evolved at 1 atm and 300 K temperature after completion of the reaction. [Assume oleum to be in excess]
\begin{qoptions}
\opt{A}{12.3 L}
\opt{B}{0.246 L}
\opt{C}{24.6 L}
\opt{D}{2.46 L}
\end{qoptions}
\end{bankq}
\begin{bankq}[id=TOAS05-P2-CHEM-S2-Q06, source={TOAS - Test 05 (03 Apr 2022)}, paper=2, qno=6, type=numerical, subtopic={Stoichiometry}, answer={30}, answer_src=key, difficulty=4, diff_parts={K1 S2 I1 C0 F1}, diff_note={concepts: gas volume contraction in reactions | steps: pyrogallol removes O2; contraction in the first reaction; contraction in the repeated reaction; composition | I: use the repeated reaction to find the unreacted gas}]
A mixture of $\mathrm{H_2}$, $\mathrm{N_2}$ and $\mathrm{O_2}$ occupying 100 ml underwent reaction so as to form $\mathrm{H_2O_2}$ (l) and $\mathrm{N_2H_2}$ (g) as the only products, causing the volume of contract by 60 ml. The remaining mixture was passed through pyragallol causing a contraction of 10 ml. To the remaining mixture excess $\mathrm{H_2}$ was added and the above reaction was repeated, causing a reduction in volume of 10 ml. (No other products are formed)\\
What is the mole percent of $\mathrm{N_2}$ in the initial mixture ?
\end{bankq}
\begin{bankq}[id=TOAS16-P1-CHEM-S2-Q05, source={TOAS - Test 16 (04 Sep 2022)}, paper=1, qno=5, type=multi, subtopic={Stoichiometry}, answer={B,D}, answer_src=key, difficulty=6, diff_parts={K2 S2 I1 C1 F1}, diff_note={concepts: iodometric titration stoichiometry, mole relations in a process | steps: iodine from hypo titration; scale to the whole sample; moles of AgI and Fe with the yields; mass of silver nitrate}]
The following process has been used to obtain iodine from oil-field brines in California :\par
\noindent NaI + AgNO$_3$ $\xrightarrow{\text{Yield}=100\%}$ AgI + NaNO$_3$\par
\noindent AgI + Fe $\xrightarrow{\text{Yield}=80\%}$ FeI$_2$ + Ag\par
\noindent FeI$_2$ + Cl$_2$ $\xrightarrow{\text{Yield}=100\%}$ FeCl$_3$ + I$_2$\par
$\dfrac{1}{2\times10^4}$ part of the sample of I$_2$ produced is dissolved in water to make 1 litre solution. 10 ml of this solution requires 20 ml 0.1N hypo solution to decolorise the blue starch iodine complex. Select the correct statement(s) :
\begin{qoptions}
\opt{A}{The mass of AgNO$_3$ initially taken is 680 g}
\opt{B}{The mass of AgNO$_3$ initially taken is 850 kg}
\opt{C}{Fe required is 2000 mole}
\opt{D}{Fe required is 2500 mole}
\end{qoptions}
\end{bankq}
\begin{bankq}[id=TOAS19-P1-CHEM-S1-Q02, source={TOAS - Test 19 (05 Feb 2023)}, paper=1, qno=2, type=multi, subtopic={Concentration terms}, answer={A,B,D}, answer_src=key, difficulty=3, diff_parts={K1 S2 I0 C0 F1}, diff_note={concepts: volume strength of hydrogen peroxide | steps: molarity from the volume strength; moles in the cylinder; oxygen used; strength left}]
One of the major use of H$_2$O$_2$ is for artificial respiration. It is stored in cylinder of 3 litre capacity completely filled with H$_2$O$_2$ solution of `45.4V'. In a hospital of kota, patient under artificial respiration takes 200 ml O$_2$ per min at 1 atm at 273 K for 2.8 hours from cylinder.

[Assume volume of solution and rate of decomposition remain constant]

Then which of the following statement(s) is/are correct.
\begin{qoptions}
\opt{A}{Initial moles of H$_2$O$_2$ in the cylinder are 12.}
\opt{B}{Volume of oxygen used for respiration is 33.6 L.}
\opt{C}{Volume strength of H$_2$O$_2$ left in solution is 11.35 V.}
\opt{D}{Volume strength of H$_2$O$_2$ left in solution is 34.05 V.}
\end{qoptions}
\end{bankq}

\section*{TAAS papers}
\begin{bankq}[id=TAAS02-P2-CHEM-S4-Q10, source={TAAS - Test 02 (30 Jul 2023)}, paper=2, qno=10, type=numerical, subtopic={Concentration terms}, answer={9}, answer_src=key, difficulty=6, diff_parts={K1 S2 I2 C1 F1}, diff_note={concepts: molarity equal to molality | steps: relation between molarity and molality with density; mole fraction; ratio of molar masses}]
The mole fraction of a solute in a solution is 0.1. At 298 K, molarity of this solution is the same as its molality. Density of this solution at 298 K is 2.0 g cm$^{-3}$. The ratio of the molecular weights of the solute and solvent, $\left(\dfrac{MW_{\text{solute}}}{MW_{\text{solvent}}}\right)$, is
\end{bankq}

\section*{TOPS papers}
\begin{bankq}[id=TOPS02-P1-CHEM-S1-Q01, source={TOPS - Test 02 (28 Aug 2022)}, paper=1, qno=1, type=single, subtopic={Stoichiometry}, answer={C}, answer_src=key, difficulty=2, diff_parts={K1 S1 I0 C1 F0}, diff_note={concepts: mole concept | steps: count hydrogens per formula unit; multiply by 0.1 mole}]
How many hydrogen atoms are present in 0.1 mole of mohr's salt ($\mathrm{FeSO_4.(NH_4)_2SO_4.6H_2O}$) ?
\begin{qoptions}
\opt{A}{$10\times N_A$}
\opt{B}{$14\times N_A$}
\opt{C}{$2\times N_A$}
\opt{D}{$0.8\,N_A$}
\end{qoptions}
\end{bankq}
\begin{bankq}[id=TOPS02-P1-CHEM-S1-Q02, source={TOPS - Test 02 (28 Aug 2022)}, paper=1, qno=2, type=single, subtopic={Stoichiometry}, answer={D}, answer_src=key, difficulty=3, diff_parts={K1 S2 I0 C1 F0}, diff_note={concepts: mole concept | steps: moles of H atoms; moles of O in water; mass}]
A sample of water contains 10 gm hydrogen atom. The weight of oxygen present in the given sample is
\begin{qoptions}
\opt{A}{160 gm}
\opt{B}{320 gm}
\opt{C}{240 gm}
\opt{D}{80 gm}
\end{qoptions}
\end{bankq}
\begin{bankq}[id=TOPS02-P1-CHEM-S1-Q03, source={TOPS - Test 02 (28 Aug 2022)}, paper=1, qno=3, type=single, subtopic={Stoichiometry}, answer={A}, answer_src=key, difficulty=2, diff_parts={K1 S1 I0 C1 F0}, diff_note={concepts: mole concept | steps: moles of water; protons per molecule}]
How many mole of protons are present in 90g of H$_2$O ?
\begin{qoptions}
\opt{A}{50 mole}
\opt{B}{60 mole}
\opt{C}{70 mole}
\opt{D}{40 mole}
\end{qoptions}
\end{bankq}
\begin{bankq}[id=TOPS02-P1-CHEM-S2-Q07, source={TOPS - Test 02 (28 Aug 2022)}, paper=1, qno=7, type=multi, subtopic={Stoichiometry}, answer={B,C}, answer_src=key, difficulty=3, diff_parts={K1 S1 I0 C1 F1}, diff_note={concepts: mole concept | steps: molar mass of P4; atoms and molecules in it}]
124 u of P$_4$ will contain
\begin{qoptions}
\opt{A}{4 $N_A$ atoms of phosphorus}
\opt{B}{4 atoms of phosphorus}
\opt{C}{1 molecule of phosphorus}
\opt{D}{$N_A$ molecules of phosphorus}
\end{qoptions}
\end{bankq}
\begin{bankq}[id=TOPS02-P1-CHEM-S2-Q08, source={TOPS - Test 02 (28 Aug 2022)}, paper=1, qno=8, type=multi, subtopic={Stoichiometry}, answer={A,C,D}, answer_src=key, difficulty=3, diff_parts={K1 S1 I0 C1 F1}, diff_note={concepts: molar volume | steps: volume of a liquid mole; molar volume of gases at STP and 298 K}]
Which is/are correct statements ?
\begin{qoptions}
\opt{A}{Molar volume of H$_2$O($\ell$) is almost equal to 18 ml at 1 atm pressure and 298 K temperature}
\opt{B}{Molar volume of N$_2$ gas is equal to 22.4 L at 1 atm pressure and 298 K temperature}
\opt{C}{Molar volume of O$_2$ gas is equal to 22.7 L at STP}
\opt{D}{Molar volume of NH$_3$ gas is equal to 22.4 L at 1 atm and 273 K.}
\end{qoptions}
\end{bankq}
\begin{bankq}[id=TOPS02-P1-CHEM-S2-Q09, source={TOPS - Test 02 (28 Aug 2022)}, paper=1, qno=9, type=multi, subtopic={Stoichiometry}, answer={A,B}, answer_src=key, difficulty=3, diff_parts={K1 S1 I0 C1 F1}, diff_note={concepts: mole concept | steps: moles of each atom in 10 moles; check each statement}]
Which is/are correct statements about 10 mol of Na$_2$SO$_4$.7H$_2$O.
\begin{qoptions}
\opt{A}{$6.022\times10^{24}$ atoms of sulphur present}
\opt{B}{$7\times6.022\times10^{24}$ molecules of H$_2$O present}
\opt{C}{176 amu of oxygen present}
\opt{D}{46 gm of sodium present}
\end{qoptions}
\end{bankq}
\begin{bankq}[id=TOPS02-P1-CHEM-S3-Q01, source={TOPS - Test 02 (28 Aug 2022)}, paper=1, qno=1, type=numerical, subtopic={Stoichiometry}, answer={4.48}, answer_src=key, difficulty=6, diff_parts={K2 S2 I1 C1 F1}, diff_note={concepts: moles, graph reading | steps: read the molar mass from the graph; moles of 5.6 g; volume at STP}]
For a gaseous substance a graph between various mole was plotted against the respective mass(g) of the substance as given below. What is the volume ($\ell$) of 5.6g of the substance at 1atm and 273K.
\bfig{115.3}{TOPS02_p1_chem_s3_q01_1.png}
\end{bankq}
\begin{bankq}[id=TOPS02-P1-CHEM-S3-Q02, source={TOPS - Test 02 (28 Aug 2022)}, paper=1, qno=2, type=numerical, subtopic={Stoichiometry}, answer={4.00}, answer_src=key, difficulty=4, diff_parts={K1 S2 I0 C1 F1}, diff_note={concepts: mole concept | steps: molar mass of adipic acid; moles in 58.4 g; hydrogen mass}]
Adipic acid HOOC(CH$_2$)$_4$COOH is used in making J nylon. Calculate the weight of hydrogen atom in gram present in 58.4 g of adipic acid.
\end{bankq}
\begin{bankq}[id=TOPS02-P1-CHEM-S3-Q03, source={TOPS - Test 02 (28 Aug 2022)}, paper=1, qno=3, type=numerical, subtopic={Stoichiometry}, answer={4.00}, answer_src=key, difficulty=3, diff_parts={K1 S1 I0 C1 F1}, diff_note={concepts: mole concept | steps: mass of iron in one mole of protein; divide by 56}]
If iron (Fe 56) is the only metal in a protein and rest are non-metals, how many iron atoms are present in one molecule of protein if the molar mass of protein is 64000 gm/mol and it contains 0.35\% iron by mass? (Atomic mass of Fe $=56$)
\end{bankq}
\begin{bankq}[id=TOPS02-P2-CHEM-S1-Q01, source={TOPS - Test 02 (28 Aug 2022)}, paper=2, qno=1, type=multi, subtopic={Stoichiometry}, answer={A,B}, answer_src=key, difficulty=3, diff_parts={K1 S1 I0 C1 F1}, diff_note={concepts: mole concept | steps: atoms of each type in the formula; ratios and total atoms}]
Select the correct statement(s) for $\mathrm{(NH_4)_3PO_4}$.
\begin{qoptions}
\opt{A}{Ratio of number of oxygen atoms to number of hydrogen atoms is 1 : 3}
\opt{B}{Ratio of number of cations to number of anions is 3 : 1}
\opt{C}{Ratio of number of gm-atoms of nitrogen to gm-atoms of oxygen is 3 : 2}
\opt{D}{Total number of atoms in one mole of $\mathrm{(NH_4)_3PO_4}$ is 20}
\end{qoptions}
\end{bankq}
\begin{bankq}[id=TOPS02-P2-CHEM-S1-Q02, source={TOPS - Test 02 (28 Aug 2022)}, paper=2, qno=2, type=multi, subtopic={Stoichiometry}, answer={A,B,C,D}, answer_src=key, difficulty=4, diff_parts={K1 S2 I0 C1 F1}, diff_note={concepts: mole concept | steps: mass of each option; molar masses; check each}]
Which of the following has a mass of 20 g. [Use $N_A=6\times10^{23}$]

[Atomic mass : Ca $=40$, S $=32$, O $=16$, He $=4$, C $=12$]
\begin{qoptions}
\opt{A}{$3\times10^{23}$ Ca$^{2+}$ ion}
\opt{B}{5.6 L SO$_3$ gas at 1 atm and 273 K}
\opt{C}{Mixture of 1 mol He gas and 1 mol CH$_4$ gas}
\opt{D}{$1.2\times10^{23}$ molecule of CaCO$_3$.}
\end{qoptions}
\end{bankq}
\begin{bankq}[id=TOPS02-P2-CHEM-S1-Q03, source={TOPS - Test 02 (28 Aug 2022)}, paper=2, qno=3, type=multi, subtopic={Stoichiometry}, answer={A,B,C,D}, answer_src=key, difficulty=4, diff_parts={K1 S2 I0 C1 F1}, diff_note={concepts: mole concept | steps: gram atoms and molecules; electrons per atom; check each statement}]
Identify the correct statements among the following
\begin{qoptions}
\opt{A}{0.5 gm-atom of oxygen contains $\dfrac{N_A}4$ oxygen molecules}
\opt{B}{Total electrons present in 0.5 g-atom of oxygen is $4N_A$}
\opt{C}{Mass of $3.01\times10^{23}$ molecules of ammonia is 8.5 g}
\opt{D}{Number of atoms present in 11.2 litres of chlorine at 1 atm \& 273K is $6.022\times10^{23}$}
\end{qoptions}
\end{bankq}
\begin{bankq}[id=TOPS02-P2-CHEM-S2-Q01, source={TOPS - Test 02 (28 Aug 2022)}, paper=2, qno=1, type=numerical, subtopic={Stoichiometry}, answer={22.00}, answer_src=key, difficulty=4, diff_parts={K1 S2 I0 C1 F1}, diff_note={concepts: mole concept | steps: moles of H in 9 g of water; moles of CO2 with the same oxygen; mass}]
Let $x$(g) be mass of CO$_2$ that must be taken so that it contains same number of O atoms as there number of H atoms in 9 g H$_2$O. Find $x$.
\end{bankq}
\begin{bankq}[id=TOPS02-P2-CHEM-S2-Q02, source={TOPS - Test 02 (28 Aug 2022)}, paper=2, qno=2, type=numerical, subtopic={Stoichiometry}, answer={21.00}, answer_src=key, difficulty=4, diff_parts={K1 S2 I0 C1 F1}, diff_note={concepts: mole concept | steps: moles of the hydrate; moles of water; number of molecules}]
$2.4\times10^y$ molecules of water are present in 252 mg of H$_2$C$_2$O$_4$.2H$_2$O then value of $y$ is-

[Given $N_A=6\times10^{23}$]
\end{bankq}
\begin{bankq}[id=TOPS02-P2-CHEM-S2-Q03, source={TOPS - Test 02 (28 Aug 2022)}, paper=2, qno=3, type=numerical, subtopic={Stoichiometry}, answer={400.00}, answer_src=key, difficulty=3, diff_parts={K1 S1 I0 C1 F1}, diff_note={concepts: percentage composition | steps: one sulphur atom is the least; molar mass equals 32 over 0.08}]
A compound possesses 8\% sulphur by mass. The least molecular mass of the compound is
\end{bankq}
\begin{bankq}[id=TOPS02-P2-CHEM-S3-Q01, source={TOPS - Test 02 (28 Aug 2022)}, paper=2, qno=1, type=numerical, subtopic={Stoichiometry}, answer={3.00}, answer_src=key, difficulty=4, diff_parts={K1 S2 I0 C1 F1}, diff_note={concepts: mole concept | steps: oxygen per mole of alum; moles needed; mass; digit sum}]
The weight (in gram) of pure potash Alum ($\mathrm{K_2SO_4.Al_2(SO_4)_3.24H_2O}$) which contains 0.64 kg oxygen. (Atomic weight of K $=39$, S $=32$, Al $=27$)

[If your answer is greater than 9 then add all the digits. Example : Ans.12 then fill $1+2=3$, Ans.123 then fill $1+2+3=6$ \& Ans. 20 then fill $2+0=2$]
\end{bankq}
\begin{bankq}[id=TOPS02-P2-CHEM-S3-Q02, source={TOPS - Test 02 (28 Aug 2022)}, paper=2, qno=2, type=numerical, subtopic={Stoichiometry}, answer={3.00}, answer_src=key, difficulty=4, diff_parts={K1 S2 I0 C1 F1}, diff_note={concepts: mole concept | steps: moles of each; oxygen atoms; ratio}]
The simplest ratio of number of oxygen atoms (O) in 16.0g of ozone (O$_3$), 28.0 g carbon monoxide (CO) and 16.0g oxygen (O$_2$) is $x:y:z$ the value of $x+y+z$ is :

(Atomic mass : C $=12$, O $=16$ and Avogadro's constant NA $=6.0\times10^{23}$ mol$^{-1}$)
\end{bankq}
\begin{bankq}[id=TOPS03-P1-CHEM-S1-Q01, source={TOPS - Test 03 (20 Nov 2022)}, paper=1, qno=1, type=single, subtopic={Stoichiometry}, answer={C}, answer_src=key, difficulty=5, diff_parts={K2 S2 I1 C1 F0}, diff_note={concepts: gaseous stoichiometry, mass percent | steps: combustion equations; volume change; solve the volumes; mass percent | mixture analysis}]
20 ml of mixture of CH$_4$ and CO$_2$ when reacted with just sufficient oxygen, total volume after the complete reaction becomes 40 ml at 1 atm and 500K (for all volume measurement). The mass percentage of CH$_4$ in the initial mixture is :
\begin{qoptions}
\opt{A}{50}
\opt{B}{40}
\opt{C}{26.66}
\opt{D}{60}
\end{qoptions}
\end{bankq}
\begin{bankq}[id=TOPS03-P1-CHEM-S1-Q02, source={TOPS - Test 03 (20 Nov 2022)}, paper=1, qno=2, type=single, subtopic={Stoichiometry}, answer={D}, answer_src=key, difficulty=5, diff_parts={K2 S2 I1 C1 F0}, diff_note={concepts: redox titration, equivalents | steps: equivalents of permanganate; equivalents of oxalate; moles of calcium; mass of CaO}]
0.518 g sample of lime-stone is dissolved in HCl and entire calcium is precipitated as CaC$_2$O$_4$. After filtering and washing the precipitate 40 ml of 0.25 N acidified KMnO$_4$ were required to titrate the precipitate in a redox titration. Calculate mass of CaO obtained from strong heating of the initial sample :
\begin{qoptions}
\opt{A}{0.13 g}
\opt{B}{0.56 g}
\opt{C}{0.26 g}
\opt{D}{0.28 g}
\end{qoptions}
\end{bankq}
\begin{bankq}[id=TOPS03-P1-CHEM-S2-Q06, source={TOPS - Test 03 (20 Nov 2022)}, paper=1, qno=6, type=multi, subtopic={Concentration terms}, answer={A,B,C,D}, answer_src=key, difficulty=9, diff_parts={K3 S3 I2 C1 F1}, diff_note={concepts: molality, neutralisation, concentration terms | steps: molality from molarity and density; moles of NaOH; sulphur trioxide and labelling; hydrogen ion concentration; check each | four statements}]
A 4 M aqueous solution of H$_2$SO$_4$ has density 0.86 g/ml. Which is/are correct for this solution :
\begin{qoptions}
\opt{A}{Molality of solution is about 8.55 m.}
\opt{B}{62.5 ml of this solution required 50 ml of 40\% w/v NaOH solution for complete neutralisation.}
\opt{C}{If 30 moles of SO$_3$ is added to 1 liter of this solution then \% labelling of this solution is 102.2\%.}
\opt{D}{If 100 ml 2 M H$_3$PO$_4$ is added to 100 ml of this solution then concentration of H$^+$ in the final solution is 7 M. (Assuming complete dissociation of acids).}
\end{qoptions}
\end{bankq}
\begin{bankpassage}[id=TOPS03-P1-CHEM-S3-G1-2, source={TOPS - Test 03 (20 Nov 2022)}, questions={1-2}]
Read the following passage and answer the questions based on it.

In a reaction, formation of maximum number of moles of any of the products depends upon the reagent which has stoichiometrically minimum number of moles among the reactants. This reagent, having stoichiometrically minimum number of moles is called Limiting Reagent', which is identified by considering number of moles of all reactants together with stoichiometry (balanced) of the reaction.
\end{bankpassage}
\begin{bankq}[id=TOPS03-P1-CHEM-S3-Q01, source={TOPS - Test 03 (20 Nov 2022)}, paper=1, qno=1, type=numerical, subtopic={Limiting reagent}, answer={0.71}, answer_src=key, difficulty=8, diff_parts={K2 S3 I2 C1 F1}, diff_note={concepts: limiting reagent, precipitation | steps: moles of each reactant; limiting reagent; moles remaining; moles of NaCl formed; total ions and concentration | precipitation reaction}, group=TOPS03-P1-CHEM-S3-G1-2]
500 mL of 0.4 M BaCl$_2$ solution was mixed with 1000 mL of 0.2 M Na$_3$PO$_4$ solution, then total concentration of ions in resultant solution is \underline{\hspace{1cm}} M? Assume Ba$_3$(PO$_4$)$_2$ completely insoluble in water.
\end{bankq}
\begin{bankq}[id=TOPS03-P1-CHEM-S3-Q02, source={TOPS - Test 03 (20 Nov 2022)}, paper=1, qno=2, type=numerical, subtopic={Limiting reagent}, answer={0.12}, answer_src=key, difficulty=6, diff_parts={K2 S2 I1 C1 F1}, diff_note={concepts: limiting reagent, stoichiometry | steps: moles of each element; balanced ratio; limiting element; moles formed}, group=TOPS03-P1-CHEM-S3-G1-2]
$3.011\times10^{23}$ atoms of Zn, 28 gm Fe and $3.011\times10^{23}$ atoms of sulphur were allowed to react in a container to form Zn(FeS$_2$)$_2$ compound. Find moles of Zn(FeS$_2$)$_2$ is formed ? (Atomic mass of Fe $=56$)
\end{bankq}
\begin{bankq}[id=TOPS03-P2-CHEM-S1-Q02, source={TOPS - Test 03 (20 Nov 2022)}, paper=2, qno=2, type=multi, subtopic={Stoichiometry}, answer={A,B}, answer_src=key, difficulty=8, diff_parts={K2 S3 I2 C1 F1}, diff_note={concepts: gaseous stoichiometry, partial pressure | steps: combustion equations; pressure of CO2; solve for the two hydrocarbons; oxygen pressure; mole fraction | pressure analysis}]
A mixture of C$_2$H$_2$ and C$_3$H$_8$ including sufficient O$_2$ occupied a certain volume at 80 mm Hg. The mixture was completely burnt to CO$_2$ and H$_2$O($\ell$) and observed pressure of CO$_2$ was found to be 42.5 mm Hg at the same temperature and volume. Select the correct statement(s), neglecting the vapour pressure of water :
\begin{qoptions}
\opt{A}{The final pressure of reaction mixture was 42.5 mmHg.}
\opt{B}{Pressure of oxygen in the initial mixture is 62.5 mmHg}
\opt{C}{Pressure of C$_2$H$_2$ in the initial mixture is 32 mmHg.}
\opt{D}{Moles fraction of C$_3$H$_8$ in the initial mixture is 0.6.}
\end{qoptions}
\end{bankq}
\begin{bankq}[id=TOPS03-P2-CHEM-S2-Q01, source={TOPS - Test 03 (20 Nov 2022)}, paper=2, qno=1, type=numerical, subtopic={Concentration terms}, answer={1.50}, answer_src=key, difficulty=5, diff_parts={K1 S2 I1 C1 F1}, diff_note={concepts: mixing of solutions | steps: moles in each; ratio for 3 M; volume}]
What would be the maximum volume (in litre) of 3M HCl solution that can be prepared by using 2M, 1L HCl and 5M, 2L HCl kept in seperate vessel.
\end{bankq}
\begin{bankq}[id=TOPS03-P2-CHEM-S3-Q01, source={TOPS - Test 03 (20 Nov 2022)}, paper=2, qno=1, type=numerical, subtopic={Concentration terms}, answer={5}, answer_src=key, difficulty=7, diff_parts={K2 S2 I2 C1 F1}, diff_note={concepts: molecular area, concentration | steps: number of molecules for the area; moles; mass; volume of solution}]
A solution of fatty acid $X$ (GMM $=252$) in benzene contain 4.2gm of acid per litre. When this solution is dropped in a water surface the benzene evaporates and fatty acid $X$ forms a mono layer film of solid type. If we wish to cover 10 m$^2$ area with monolayer film of fatty acid $X$, what volume of solution (in ml) should we use. Each fatty acid molecule occupy 0.2 nm$^2$ area. $N_A=6\times10^{23}$.
\end{bankq}
\begin{bankq}[id=TOPS03-P2-CHEM-S3-Q02, source={TOPS - Test 03 (20 Nov 2022)}, paper=2, qno=2, type=numerical, subtopic={Concentration terms}, answer={2}, answer_src=key, difficulty=6, diff_parts={K2 S2 I1 C1 F1}, diff_note={concepts: basicity of an acid, titration | steps: moles of acid; moles of alkali; ratio; basicity}]
The formula weight of an acid is 82. In a titration, 100 cm$^3$ of a solution of this acid containing 39 g of the acid per litre were completely neutralised by 95.0 cm$^3$ of aqueous NaOH containg 40.0g of NaOH per litre. Then the basicity of the acid is \underline{\hspace{1cm}}.

(Answer to nearest integer)
\end{bankq}
\begin{bankpassage}[id=TOPS04-P1-CHEM-S2-G7-8, source={TOPS - Test 04 (18 Dec 2022)}, questions={7-8}]
Potassium dichromate ($\mathrm{K_2Cr_2O_7}$) is an orange coloured compound, very frequently used in laboratory as an oxidising agent as well as in a redox titration. It is generally prepared from chromite ($\mathrm{FeCr_2O_4}$) ore according to the following reactions:

(1) Fusion of chromite ore with sodium carbonate in excess of air.

$\mathrm{FeCr_2O_4+Na_2CO_3+O_2\longrightarrow Na_2CrO_4+Fe_2O_3+CO_2}$

(2) Acidifying filtered sodium chromate solution with sulphuric acid.

$\mathrm{Na_2CrO_4+H_2SO_4\longrightarrow Na_2Cr_2O_7+Na_2SO_4+H_2O}$

(3) Treating sodium dichromate with potassium chloride.

$\mathrm{Na_2Cr_2O_7+KCl\longrightarrow K_2Cr_2O_7+NaCl}$

Answer the following questions using above information. [Fe = 56, Cr = 52]
\end{bankpassage}
\begin{bankq}[id=TOPS04-P1-CHEM-S2-Q07, source={TOPS - Test 04 (18 Dec 2022)}, paper=1, qno=7, type=single, subtopic={Limiting reagent}, answer={A}, answer_src=key, difficulty=6, diff_parts={K2 S3 I1 C1 F0}, diff_note={concepts: stoichiometry, limiting reagent | steps: balance the fusion reaction; moles of each reactant; limiting reagent; oxygen needed; air volume}, group=TOPS04-P1-CHEM-S2-G7-8]
If you are initially provided with 224 gm of pure chromite ore \& 169.6 gm of sodium carbonate, the minimum volume of air required at 1 atm and 273 K to consume at least one of the reactant completely, if air contains 20\% by volume of oxygen gas is
\begin{qoptions}
\opt{A}{156.8 L}
\opt{B}{31.36 L}
\opt{C}{28 L}
\opt{D}{152.4 L}
\end{qoptions}
\end{bankq}
\begin{bankq}[id=TOPS04-P1-CHEM-S2-Q08, source={TOPS - Test 04 (18 Dec 2022)}, paper=1, qno=8, type=single, subtopic={Limiting reagent}, answer={A}, answer_src=key, difficulty=6, diff_parts={K2 S2 I2 C1 F0}, diff_note={concepts: limiting reagent, multi-step synthesis | steps: moles through each step; limiting reagent at each step; K2Cr2O7 yield}, group=TOPS04-P1-CHEM-S2-G7-8]
If the number of moles of reactants available for reaction are: \{$\mathrm{FeCr_2O_4}$ = 0.25 moles, $\mathrm{O_2}$ = 0.35 moles, $\mathrm{Na_2CO_3}$ = 0.60 moles, $\mathrm{H_2SO_4}$ = 0.2 moles, KCl = 0.1 moles\}, then the maximum number of moles of $\mathrm{K_2Cr_2O_7}$ that can be produced is
\begin{qoptions}
\opt{A}{0.05 moles}
\opt{B}{0.1 moles}
\opt{C}{0.2 moles}
\opt{D}{0.5 moles}
\end{qoptions}
\end{bankq}
\begin{bankq}[id=TOPS04-P2-CHEM-S1-Q02, source={TOPS - Test 04 (18 Dec 2022)}, paper=2, qno=2, type=single, subtopic={Concentration terms}, answer={D}, answer_src=key, difficulty=5, diff_parts={K1 S2 I2 C1 F0}, diff_note={concepts: molarity and molality | steps: molality and molarity relation; moles of chloride; density}]
An aqueous solution contains 2M-$\mathrm{AlCl_3}$ and 3m-$\mathrm{SnCl_4}$. If the molarity and molality of $\mathrm{Cl^-}$ ion in the solution is same, the density of solution is (Al = 27, Sn = 119)
\begin{qoptions}
\opt{A}{1.05 gm/ ml}
\opt{B}{1.15 gm/ ml}
\opt{C}{1.95 gm / ml}
\opt{D}{2.05 gm/ ml}
\end{qoptions}
\end{bankq}
\begin{bankq}[id=TOPS04-P2-CHEM-S1-Q04, source={TOPS - Test 04 (18 Dec 2022)}, paper=2, qno=4, type=single, subtopic={Concentration terms}, answer={B}, answer_src=key, difficulty=4, diff_parts={K1 S2 I1 C1 F0}, diff_note={concepts: hardness removal | steps: moles of calcium bicarbonate; moles of lime; mass}]
What is the mass of $\mathrm{Ca(OH)_2}$ required for 10 litre of water to remove temporary hardness of 100 ppm due to $\mathrm{Ca(HCO_3)_2}$ ?
\begin{qoptions}
\opt{A}{1.62 g}
\opt{B}{0.74 g}
\opt{C}{7.4 g}
\opt{D}{0.46 g}
\end{qoptions}
\end{bankq}
\begin{bankq}[id=TOPS04-P2-CHEM-S1-Q06, source={TOPS - Test 04 (18 Dec 2022)}, paper=2, qno=6, type=single, subtopic={Stoichiometry}, answer={C}, answer_src=key, difficulty=6, diff_parts={K2 S2 I2 C1 F0}, diff_note={concepts: stoichiometry, mass percent | steps: equal moles of two oxides; oxygen and phosphorus balance; mass; percent}]
A 4g mixture of $\mathrm{P_4}$ and $\mathrm{O_2}$ react completely (No reactant left unreacted) and form equal moles of $\mathrm{P_4O_6}$ and $\mathrm{P_4O_{10}}$. Find mass\% of $\mathrm{P_4}$ in original sample. (P = 31)
\begin{qoptions}
\opt{A}{60.2\%}
\opt{B}{51.4\%}
\opt{C}{49.2\%}
\opt{D}{30.4\%}
\end{qoptions}
\end{bankq}
\begin{bankq}[id=TOPS04-P2-CHEM-S2-Q01, source={TOPS - Test 04 (18 Dec 2022)}, paper=2, qno=1, type=numerical, subtopic={Concentration terms}, answer={1170}, answer_src=solved, working={5 m means 5 mol $\mathrm{H_2O_2}$ per kg water. 2.5 mol of $\mathrm{O_2}$ needs 5 mol $\mathrm{H_2O_2}$, i.e. 1000 g water plus 170 g $\mathrm{H_2O_2}$, so 1170 g of solution}, difficulty=7, diff_parts={K2 S2 I2 C1 F1}, diff_note={concepts: molality, decomposition | steps: moles of peroxide needed; mass per kilogram water; solution mass}]
Beautician of the Queen Elizabeth had an $\mathrm{H_2O_2}$ solution of molality `5m' \& density 1.17 gm/ml. If 2.5 mol of $\mathrm{O_2}$ is needed for bleaching of the Queen's hair, then find the minimum mass (in grams) of the $\mathrm{H_2O_2}$ solution the beautician needs for the bleaching purpose. Do not forget that the $\mathrm{H_2O_2}$ decomposes to give water and oxygen gas.
\end{bankq}
\begin{bankq}[id=TOPS05-P1-CHEM-S1-Q01, source={TOPS - Test 05 (14 Jan 2023)}, paper=1, qno=1, type=multi, subtopic={Stoichiometry}, answer={B,C}, answer_src=key, difficulty=7, diff_parts={K3 S2 I1 C1 F1}, diff_note={concepts: equivalents, volume strength, stoichiometry | steps: equivalent of the acid; volume strength; combustion; moles of oxygen atoms}]
The correct option(s) is/are :
\begin{qoptions}
\opt{A}{In 1 equivalent of $\mathrm{H_2SO_4}$, 2 mole of H-atom is present.}
\opt{B}{Volume strength of $\dfrac15$M - $\mathrm{H_2O_2}$ solution is `2.27 V'}
\opt{C}{14 gm of ethene required, 48 gm of oxygen gas for complete combustion.}
\opt{D}{Number of moles of oxygen atom in 276 gm of $\mathrm{N_2O_4}$ is 6.}
\end{qoptions}
\end{bankq}
\begin{bankq}[id=TOPS05-P2-CHEM-S1-Q01, source={TOPS - Test 05 (14 Jan 2023)}, paper=2, qno=1, type=single, subtopic={Concentration terms}, answer={C}, answer_src=key, difficulty=5, diff_parts={K1 S2 I2 C1 F0}, diff_note={concepts: molality mixing | steps: moles per kg of solvent; solvent ratio; solution masses; compare}]
A 1.5 m solution of acetic acid (I) is mixed with 3 m solution of the acetic acid (II) to prepare 2 m solution ($m$ is molality of solution). Select the correct statement(s) :

(I) Mass ratio of solvents mixed $\left(\dfrac{\mathrm I}{\mathrm{II}}\right)$ is $\dfrac12$

(II) Mass ratio of solvents mixed $\left(\dfrac{\mathrm I}{\mathrm{II}}\right)$ is $\dfrac21$

(III) Mass ratio of solutions mixed $\left(\dfrac{\mathrm I}{\mathrm{II}}\right)$ is $\dfrac{109}{59}$

(IV) Mass ratio of solutions mixed $\left(\dfrac{\mathrm I}{\mathrm{II}}\right)$ is $\dfrac{59}{109}$
\begin{qoptions}
\opt{A}{I, III}
\opt{B}{I, IV}
\opt{C}{II, III}
\opt{D}{II, IV}
\end{qoptions}
\end{bankq}

\section*{TEAS papers}
\begin{bankq}[id=TEAS01-P1-CHEM-S1-Q05, source={TEAS - Test 01 (16 Apr 2023)}, paper=1, qno=5, type=multi, subtopic={Stoichiometry}, answer={A,C}, answer_src=key, difficulty=5, diff_parts={K1 S2 I1 C1 F1}, diff_note={concepts: redox stoichiometry | steps: moles of sulphur dioxide; permanganate moles; barium chloride moles; check each}]
8.0 gm of solid sulphur is first oxidised to $\mathrm{SO_2}$ and then it is divided into two equal parts. One part is sufficient for complete decolorisation of 200 ml of acidified $\mathrm{KMnO_4}$ solution. Another part is oxidised to $\mathrm{SO_3}$ and the $\mathrm{SO_3}$ formed is sufficient for complete precipitation of all $\mathrm{BaCl_2}$ present in 100 ml solution as $\mathrm{BaSO_4}$. Which of the following statements is/are correct ? ($\mathrm S=32$, $\mathrm{Ba}=138$)
\begin{qoptions}
\opt{A}{The molarity of $\mathrm{KMnO_4}$ solution is 0.25}
\opt{B}{The molarity of $\mathrm{BaCl_2}$ solution is 0.25}
\opt{C}{The weight of $\mathrm{BaSO_4}$ precipitated out is 29.25 gm.}
\opt{D}{The same moles of $\mathrm{KMnO_4}$ and $\mathrm{BaCl_2}$ are reacted.}
\end{qoptions}
\end{bankq}
\end{document}
